% Article VII. Incorporation of pre-existing results, 10-09-2026.
% Complete source owners preserved in fuentes_conservadas:
% 17_pantalla_cubica_soldadura.tex (P), 18_clausura_energetica.tex (E).
% P: incidence, quotient, cellular soldering, and refinement.
% E: response, uniqueness, minimal extension, and tracial expectation.
% Earlier dependencies of VII: APP--TRIT--TPK kernel; cubic realization
% of the route, transport cocycle, and compression q_vis=5/9.
% This fragment does not by itself constitute the self-contained article.
% The screen identity is distinguished from the matter spin contraction.

\section{Boundary soldering and energy response}
\label{vii:sec:pantalla-respuesta}

The cubic realization of transport preserves six oriented incidences.
Its geometric quotient, Lorentzian form, and energy response
are obtained from the same incidence matrix. Transport originates in
the enriched APP--TRIT--TPK route: face, sign, carry, and
memory accompany each edge. This section develops the
linear operation acting on that boundary representation.

\subsection{Cubic incidence and the rank-four quotient}
\label{vii:sec:cociente-pantalla}

Let \(H_F=\mathbb R^F\), with
\(F=\{(i,\epsilon):1\leq i\leq3,\ \epsilon\in\{-1,1\}\}\),
its Euclidean inner product, and its basis \(e_{i,\epsilon}\).
Face opposition is the orthogonal involution
\(Re_{i,\epsilon}=e_{i,-\epsilon}\). Define
\[
 u=\sum_{i,\epsilon}e_{i,\epsilon},\qquad
 P_u=\frac{uu^{\mathsf T}}6,\qquad P_-=\frac{I-R}{2},\qquad
 P_0=\frac{I+R}{2}-P_u,\qquad P_\square=P_u+P_-.
\]
The vertices are \(V=\{-1,1\}^3\). The incidence matrix
\(M:\mathbb R^V\longrightarrow H_F\) is determined by
\[
 M_{(i,\epsilon),\nu}=\mathbf1_{\{\nu_i=\epsilon\}}.
\]

\begin{proposition}[Incidence decomposition]
\label{vii:prop:incidencia}
We have
\[
 MM^{\mathsf T}=12P_u+4P_-,
 \qquad
 \ker M^{\mathsf T}=P_0H_F.
\]
Consequently, the quotient
\(Q=H_F/P_0H_F\), orthogonally identified with
\(P_\square H_F\), has dimension four.
\end{proposition}

\begin{proof}
A face contains four vertices. Two opposite faces have no common
vertices, and two faces on different axes have two common vertices.
Writing \(\mathbf J=uu^{\mathsf T}\), we obtain
\[
 MM^{\mathsf T}=4I+2(\mathbf J-I-R)
 =2(I-R)+2\mathbf J=4P_-+12P_u.
\]
The three projectors \(P_u,P_-,P_0\) are orthogonal and sum to \(I\);
their ranks are respectively one, three, and two. The identity
\(\|M^{\mathsf T}x\|^2=\langle x,MM^{\mathsf T}x\rangle\)
identifies the stated kernel. The complementary rank is four.
\end{proof}

On \(Q\), the form
\[
 B=P_--P_u
\]
has signature \((3,1)\). The time orientation is fixed by \(u\).
With
\[
 t=\frac{u}{\sqrt6},\qquad
 d_i=\frac{e_{i,+}-e_{i,-}}{\sqrt2},\qquad
 W=\sqrt6P_u+\sqrt2P_-,
\]
we have \(We_{i,\epsilon}=t+\epsilon d_i\). These six vectors
are null for \(B\), whereas
\[
 3t+\nu_1d_1+\nu_2d_2+\nu_3d_3
\]
has quadratic norm \(-6\). The same incidence satisfies
\(MM^{\mathsf T}=2W^*W\) on \(Q\). Proper rotations of the cube
commute with the projectors, preserve both forms, and fix \(t\).
These identities determine the metric of the cubic representation.

\subsection{Cellular soldering and memory transport}
\label{vii:sec:soldadura-celular}

For each oriented edge \(a\), representation of the route provides
the face \(\lambda(\underline a)\), the sign
\(\sigma(a)\), the source and target fibres, and the transport
\(U(a):Q_{s(a)}\longrightarrow Q_{t(a)}\).
Inversion satisfies \(U(\bar a)=U(a)^{-1}\).
Boundary half-edges retain their formal inversion and their incidence
on the boundary. With \(\ell_m=\ell_0\,9^{-m}>0\), the soldering is
\begin{equation}
 \beta_m(a)=\ell_m\sigma_m(a)
              n_{\lambda(\underline a)}^{s(a)},\qquad
 n_{i,\epsilon}=t+\epsilon d_i.
 \label{vii:eq:soldadura-arista}
\end{equation}
The face and sign are transported together with the route memory.
The inversion convention requires
\begin{equation}
 \beta_m(\bar a)=-U_m(a)\beta_m(a).
 \label{vii:eq:inversion-soldadura}
\end{equation}
Applying this relation twice recovers \(\beta_m(a)\);
the carry register separately preserves the ordered composition.

\begin{proposition}[Naturality of soldering under refinement]
\label{vii:prop:refinamiento-soldadura}
Let \(a=a_1\cdots a_9\) be a compatible refinement of an edge.
Suppose that the nine successor incidences, transported to the predecessor
fibre, preserve face and sign, and that
\(\ell_{m+1}=\ell_m/9\). Let
\(\mathcal U_j:Q_{m,s(a)}\longrightarrow Q_{m+1,s(a_j)}\)
be the accumulated transport including the refinement identification
and the path to the source of successor \(j\). Then
\[
 \beta_m(a)=
 \sum_{j=1}^{9}\mathcal U_j^{-1}\beta_{m+1}(a_j).
\]
The equality is compatible with inversion and with concatenation
of the memory register.
\end{proposition}

\begin{proof}
Each summand, expressed in the fibre at the predecessor's source, is
\((\ell_m/9)\sigma_m(a)n_{\lambda(\underline a)}\).
The sum of nine terms reproduces the predecessor's soldering.
Inversion reverses the order of the transports and replaces each
by its inverse; relation \eqref{vii:eq:inversion-soldadura} supplies
the overall sign. Memory is concatenated in that same order, so
refinement preserves both the vector representation and its
enriched register.
\end{proof}

\subsection{Positive response, uniqueness, and minimal extension}
\label{vii:sec:respuesta-positiva}

The linear screen soldering is
\(\bar\beta_m=\ell_mW|_Q\). It is invertible, with
\[
 \bar\beta_m^{-1}
 =\ell_m^{-1}\left(\frac1{\sqrt6}P_u+\frac1{\sqrt2}P_-\right).
\]
The eigenvalue \(5/9\) of the antisymmetric sector of the normalized block
\(B_c/9\), constructed in Section~\ref{vii:sec:bloque-compresion},
determines the compression
coefficient \(q_{\mathrm{vis}}=5/9\). We reserve this notation to
distinguish it from the Euclidean quotient \(q_9(n)\) of the arithmetic kernel.
Its realization on the screen uses an isometry
\(V_9:Q\longrightarrow\mathcal K\). The visible contribution
\(C_9=q_{\mathrm{vis}}V_9\) and its complement
\(D_9=\sqrt{1-q_{\mathrm{vis}}^2}\,V_9\) satisfy
\[
 C_9^*C_9+D_9^*D_9=I_Q,\qquad D_9^*D_9=\frac{56}{81}I_Q.
\]
The coefficient \(56\) is \(9^2-5^2\); it belongs to the quadratic defect
of this compression reader.

\begin{theorem}[Energy response determined by soldering]
\label{vii:thm:respuesta-unica}
There exists a unique operator \(\Lambda_m:Q\longrightarrow Q\) such that
\(\bar\beta_m^*\Lambda_m\bar\beta_m=D_9^*D_9\). It is positive definite, and
\begin{equation}
 \Lambda_m=\frac{28}{243\ell_m^2}(P_u+3P_-)
 =\frac{112}{81\ell_m^2}
       (MM^{\mathsf T}|_Q)^{-1}.
 \label{vii:eq:lambda-pantalla}
\end{equation}
On \(H_F\), the corresponding identity has right-hand side
\((56/81)P_\square\).
\end{theorem}

\begin{proof}
Congruence by the invertible soldering necessarily determines
\[
 \Lambda_m=\bar\beta_m^{-*}\frac{56}{81}I_Q\bar\beta_m^{-1}.
\]
Substituting the inverse, the coefficient of \(P_u\) is
\(56/(486\ell_m^2)=28/(243\ell_m^2)\); that of \(P_-\) is
\(28/(81\ell_m^2)\). Both are positive. By
Proposition \ref{vii:prop:incidencia},
\((MM^{\mathsf T}|_Q)^{-1}=P_u/12+P_-/4\),
which gives the second expression. Extension to six faces annihilates
\(P_0H_F\); the identity therefore extends through \(P_\square\).
\end{proof}

\begin{proposition}[Characterization by minimal extension]
\label{vii:prop:minima-extension}
For every \(b\in Q\),
\[
 \langle b,\Lambda_m b\rangle
 =\frac{112}{81\ell_m^2}
   \min_{Mx=b}\|x\|^2.
\]
The minimizing extension is unique:
\(x_{\min}=M^{\mathsf T}(MM^{\mathsf T}|_Q)^{-1}b\).
\end{proposition}

\begin{proof}
The stated vector satisfies \(Mx_{\min}=b\) and belongs to
\(\operatorname{im}M^{\mathsf T}=(\ker M)^\perp\).
Every other solution is \(x_{\min}+z\), with \(z\in\ker M\); by
orthogonality, its squared norm is \(\|x_{\min}\|^2+\|z\|^2\).
Moreover,
\[
 \|x_{\min}\|^2
 =\langle b,(MM^{\mathsf T}|_Q)^{-1}b\rangle.
\]
The formula for \(\Lambda_m\) completes the proof.
\end{proof}

The operator \(MM^{\mathsf T}|_Q\) realizes the Dirichlet--to--Neumann
response of this incidence formulation; its inverse
is the Neumann--to--Dirichlet response. The preceding energy corresponds
to the latter, with the normalization of soldering and of the compression
defect. The Schur complement of
\(\left(\begin{smallmatrix}I&M^{\mathsf T}\\M&0\end{smallmatrix}\right)\)
is \(-MM^{\mathsf T}\), providing the same variational elimination.

\subsection{Conservation of density under refinement}
\label{vii:sec:densidad-refinamiento}

Consider the sector of cubic transports, orthogonal for the energy
inner product and compatible with the form \(B\). Fix a
common label space \(Q^{\mathrm{lab}}\), identified
isometrically with the predecessor's screen. The solderings have
types
\(\bar\beta_m:Q^{\mathrm{lab}}\longrightarrow Q_m\) and
\(\bar\beta_{m+1,j}:Q^{\mathrm{lab}}\longrightarrow Q_{m+1,j}\);
the operators \(\Lambda\) act in their respective target fibres.
If
\(U_j:Q_j\longrightarrow Q\) transports the successor fibre to the predecessor fibre,
then
\[
 \bar\beta_{m+1,j}=\frac19U_j^{-1}\bar\beta_m,\qquad
 \Lambda_{m+1,j}=81U_j^{-1}\Lambda_mU_j.
\]
The factor \(81\) comes from \(\ell_{m+1}^{-2}=81\ell_m^{-2}\).

\begin{theorem}[Invariance of the normalized density]
\label{vii:thm:densidad}
Each successor preserves the predecessor's quadratic response:
\[
 \bar\beta_{m+1,j}^*\Lambda_{m+1,j}\bar\beta_{m+1,j}
 =\frac{56}{81}I_{Q^{\mathrm{lab}}}.
\]
For \(r\) refinements, the joint response, expressed in the predecessor
fibres, is
\[
 E_{m+r}=\frac{56}{81}I_{Q^{\mathrm{lab}}}\otimes I_{9^r}.
\]
The expectation with normalized trace
\(\tau_{9^r}(A)=9^{-r}\operatorname{Tr}(A)\) satisfies
\[
 (\operatorname{id}\otimes\tau_{9^r})(E_{m+r})
 =\frac{56}{81}I_{Q^{\mathrm{lab}}}
\]
at every finite depth.
\end{theorem}

\begin{proof}
A successor's congruence contains the factors \(1/9,81,1/9\);
their product is one. Orthogonality cancels the intermediate
transports. The direct sum of nine equal responses identifies
the first refinement with \((56/81)I_Q\otimes I_9\).
Iteration produces the tensor with \(I_{9^r}\).
Finally, \(\tau_{9^r}(I_{9^r})=1\). Expectations compose
through the identity
\(\tau_{9^{r+s}}=\tau_{9^r}\otimes\tau_{9^s}\).
\end{proof}

The unital inclusions \(A\mapsto A\otimes I_9\) are compatible with
this collection of traces. In the inductive limit of the matrix algebras,
the response is the central element
\((56/81)I_Q\otimes\mathbf1\). The unnormalized extensive sum instead gives
\(9^r(56/81)I_Q\). This distinction specifies the quantity
preserved: the quadratic screen density.

\subsection{Variational realization of the response}
\label{vii:sec:realizacion-respuesta}

The matter realization uses an isometry
\(J_{\mathrm{scr}}:Q\longrightarrow\mathcal H_\partial\)
compatible with the spin action and refinement.
The boundary co-tetrad is \(e=J_{\mathrm{scr}}\bar\beta_m\).
Let \(\Lambda_{\mathrm{CE}}\) be the boundary response obtained
by second variation of the matter and geometric action, with its
boundary conditions and elimination of the interior variables.
The compatibility defect is
\begin{equation}
 \mathcal R_{\mathrm{grav}}
 =\bar\beta_m^*
   (J_{\mathrm{scr}}^*\Lambda_{\mathrm{CE}}J_{\mathrm{scr}}-\Lambda_m)
   \bar\beta_m
 =e^*\Lambda_{\mathrm{CE}}e-D_9^*D_9.
 \label{vii:eq:defecto-respuesta}
\end{equation}
By invertibility of \(\bar\beta_m\),
\[
 \mathcal R_{\mathrm{grav}}=0
 \quad\Longleftrightarrow\quad
 J_{\mathrm{scr}}^*\Lambda_{\mathrm{CE}}J_{\mathrm{scr}}=\Lambda_m.
\]
The isometry \(J_{\mathrm{scr}}\) transports the boundary states;
the variational spin current is a bivector-valued three-form.
Each enters the realization of the response with its own
domain.

The construction determines a positive energy form, its minimizing
interior extension, and the density conserved under refinement.
The amplitude of a specific matter contribution is obtained by evaluating
the current of that realization in its variational quadratic form.
This distinguishes the operator fixing the response, the current
on which it acts, and the gravitational observable expressed by the evaluation.
